\chapter{Practice: Seven Projects}
\label{ch:examples}

In this chapter we will build, step by step, seven circuits, from simple logic
gates to a finite-state machine and PWM (pulse-width modulation: it is
covered in project~6). Each one relies on blocks from the first
part of the book. Each project lives in its own folder
inside \code{code/}: a \code{.v} file with the circuit description and a \code{.cst} file with
the pins (how it is organized is explained in Section~\ref{sec:cst}).

\begin{figure}[H]
\centering
\begin{tikzpicture}
\begin{axis}[
  ybar stacked, bar width=9mm, width=\textwidth, height=6.2cm,
  symbolic x coords={Logic, Blinker, {Running light}, Button, {Traffic light}, PWM, MUX},
  xtick=data, ylabel={Number of cells},
  legend style={at={(0.5,1.03)}, anchor=south, legend columns=3, font=\sffamily\small, draw=none},
  label style={font=\sffamily\small}, tick label style={font=\sffamily\small},
  axis lines*=left, ymin=0, grid=major, grid style={FpgaGray!20},
]
  \addplot[fill=FpgaBlue!70, draw=FpgaBlue] coordinates {(Logic,5) (Blinker,114) ({Running light},60) (Button,21) ({Traffic light},32) (PWM,70) (MUX,8)};
  \addplot[fill=FpgaGreen!60, draw=FpgaGreen] coordinates {(Logic,0) (Blinker,26) ({Running light},29) (Button,31) ({Traffic light},30) (PWM,34) (MUX,25)};
  \addplot[fill=FpgaOrange!60, draw=FpgaOrange] coordinates {(Logic,0) (Blinker,25) ({Running light},22) (Button,25) ({Traffic light},35) (PWM,49) (MUX,25)};
  \legend{LUT1--LUT4, flip-flops, ALU cells (carry)}
\end{axis}
\end{tikzpicture}
\caption{How many resources each project uses (according to the Yosys synthesis report).
Even the ``heaviest'' example uses less than 1\,\% of the FPGA}
\label{fig:usage}
\end{figure}

% =============================================================================
\section{Project 1. Logic Gates on Buttons}

\textbf{Goal:} to see that the gates AND, OR, NOT, NAND, NOR and
XOR really can be ``built'' with code. Buttons S1 and S2 will be
the inputs $a$ and $b$, and six LEDs will be the outputs of the six functions.

\begin{figure}[htbp]
\centering
\begin{tikzpicture}
  \coordinate (A) at (0,0); \coordinate (B) at (0,-0.4);
  \node[left, font=\sffamily\small] at (-0.1,-3.0) {\begin{tabular}{r}S1 $\to a$\\S2 $\to b$\end{tabular}};
  \draw[very thick, FpgaBlue] (0,0.7) -- (0,-6.6) node[below, lbl] {$a$};
  \draw[very thick, FpgaRed]  (0.4,0.7) -- (0.4,-6.6) node[below, lbl] {$b$};
  \foreach \i/\g/\t in {0/gAND/AND, 1/gOR/OR, 2/gNOT/NOT, 3/gNAND/NAND, 4/gNOR/NOR, 5/gXOR/XOR} {
    \pgfmathsetmacro{\y}{-\i*1.25}
    \ifnum\i=2
      \node[gNOT] (g\i) at (2.6,\y) {};
      \draw[wire] (g\i.input) -- (0,\y); \fill[FpgaBlue] (0,\y) circle (1.5pt);
    \else
      \node[\g] (g\i) at (2.6,\y) {};
      \draw[wire] (g\i.input 1) -- (0,0 |- g\i.input 1); \fill[FpgaBlue] (0,0 |- g\i.input 1) circle (1.5pt);
      \draw[wire] (g\i.input 2) -- (0.4,0 |- g\i.input 2); \fill[FpgaRed] (0.4,0 |- g\i.input 2) circle (1.5pt);
    \fi
    \node[gNOT] (n\i) at (4.6,\y) {};
    \draw[wire] (g\i.output) -- (n\i.input);
    \draw[wire] (n\i.output) -- (5.6,\y) node[right, font=\sffamily\small] {led[\i]\quad (\t)};
  }
  \node[lbl, align=center] at (4.6,0.8) {inversion\\for the LED};
\end{tikzpicture}
\caption{The circuit of project 1}
\end{figure}

\vfile{\codepath/01_gates/gates.v}{code/01\_gates/gates.v}

Load this file into Logra and click ``Synthesize''. On the RTL tab the circuit
is broken down into individual gates: XOR, OR and AND, and each of the six LEDs
gets its own inverter (Fig.~\ref{fig:ex-gates-rtl}). This is how the six-bit
\code{led} vector turns into six independent wires. Note the
double inverters: the NAND and NOR functions are written with \code{\textasciitilde},
and one more inversion is added because the LEDs are active low.

\shot[0.5\textwidth]{ex-gates-rtl}{Project~1 in Logra: one gate for each bit of the result}{fig:ex-gates-rtl}

The pin file. There is no clock signal here: the circuit is purely combinational, and
the file describes only the buttons and the LEDs. Note
\code{PULL\_MODE=DOWN} on the buttons: without it, a released button would give a
random level, and the LEDs would flicker.

\vfile[c]{\codepath/01_gates/gates.cst}{code/01\_gates/gates.cst}

\begin{table}[H]
\centering
\caption{Which LEDs should be on ($\bullet$ means on)}
\label{tab:gates}
\begin{tabular}{cc|cccccc}
\toprule
S1 ($a$) & S2 ($b$) & led0 AND & led1 OR & led2 NOT $a$ & led3 NAND & led4 NOR & led5 XOR \\
\midrule
0 & 0 &           &           & $\bullet$ & $\bullet$ & $\bullet$ &           \\
0 & 1 &           & $\bullet$ & $\bullet$ & $\bullet$ &           & $\bullet$ \\
1 & 0 &           & $\bullet$ &           & $\bullet$ &           & $\bullet$ \\
1 & 1 & $\bullet$ & $\bullet$ &           &           &           &           \\
\bottomrule
\end{tabular}
\end{table}

\begin{infobox}[What we got after synthesis]
Yosys reports that the circuit took only \textbf{5 LUT2}, not 6 gates
plus 6 inverters. The synthesizer merged each function with the inversion at its output into
a single LUT (for example, $\overline{a\cdot b}$ became one table). And for
\code{led[2]} the result was $\overline{\overline{a}} = a$---the LED is connected
to the button simply by a wire, and no LUT was needed at all!
\end{infobox}

\begin{taskbox}
\begin{enumerate}
  \item Load the project and check all the rows of
        Table~\ref{tab:gates} one by one.
  \item Replace the function on \code{led[5]} with the majority function of three
        signals: $a$, $b$ and $a \oplus b$. What do you get, and why?
\end{enumerate}
\end{taskbox}

\begin{diybox}[Security on two buttons]
\textbf{Task.} Build project 1 and check all four rows of Table~\ref{tab:gates}. Then change the functions so that S1 is the ``door'', S2 is the ``guard'', and an LED shows the alarm.

\textbf{In Logra.} The ``From scratch'' card, Code tab: enter \code{gates.v}; the Pins tab holds the assignment of buttons and LEDs (the \code{gates.cst} file from the book). For the alarm, keep the AND function on \code{led[0]}: the board LEDs light with a zero, so the inverter in front of the pin is already in the project.

\textbf{Expected result.} The alarm LED is on only when both buttons are pressed; the other five functions match the table.

\textbf{How to check.} Waveforms: ``Sweep inputs'' gives four rows; compare them with the table. RTL: one gate for each bit. Gowin: how many LUTs are used (about five according to the report).

\stars{1} repeat: check all 4 rows of the table. \stars{2} change: replace \code{led[5]} with the majority function of $a$, $b$ and $a\oplus b$ and explain the result (you get OR). \stars{3} project: the ``Security Panel'' stage---\code{alarm = door \& (night | armed)} on buttons S1, S2 and a third input; acceptance criterion: the eight combinations match the truth table from Chapter~\ref{ch:comb}.
\end{diybox}

% =============================================================================
\section{Project 2. A Blinking LED}

\textbf{Goal:} to get to know the clock signal and the counter. This is the
``Hello, world!'' of the FPGA world.

The oscillator produces \SI{27}{MHz}, that is, $27\,000\,000$ pulses per second.
The human eye cannot see blinking that fast. We need to \term{divide the frequency}:
we count clock cycles and every $13\,500\,000$ cycles (0.5\,s) we toggle the
LED. How many bits wide must the counter be?
\[
  2^{23} = 8\,388\,608 < 13\,500\,000 < 2^{24} = 16\,777\,216
\]
\[
  \Rightarrow\quad \text{24 bits are enough (we took 25 to have some margin).}
\]

\begin{figure}[H]
\centering
\begin{tikzpicture}[node distance=10mm]
  \node[hblock] (clk) {Oscillator\\27 MHz};
  \node[block, right=of clk, minimum width=30mm] (cnt) {Counter\\\footnotesize 0 \ldots 13\,499\,999};
  \node[block, right=of cnt] (cmp) {Compare\\\footnotesize cnt == HALF$-$1};
  \node[gblock, right=of cmp] (tff) {T flip-flop\\\footnotesize state};
  \node[oblock, below=8mm of tff] (led) {6 LEDs};
  \draw[arr] (clk) -- (cnt); \draw[bus, -{Stealth}] (cnt) -- (cmp); \draw[arr] (cmp) -- (tff);
  \draw[arr] (tff) -- node[right, lbl] {\textasciitilde state} (led);
  \draw[arr, dashed] (cmp.south) |- node[pos=0.75, below, lbl] {reset to 0} ($(cnt.south)+(0,-0.6)$) -- (cnt.south);
\end{tikzpicture}
\caption{Structure of the frequency divider}
\end{figure}

\vfile{\codepath/02_blink/blink.v}{code/02\_blink/blink.v}

The pin file: the clock oscillator on pin 4 and six LEDs. The buttons
are not used in this project, so they are not in the file---the router would
skip extra lines, but a tidy project should not have them.

\vfile[c]{\codepath/02_blink/blink.cst}{code/02\_blink/blink.cst}

\begin{figure}[H]
\centering
\begin{tikztimingtable}[timing/slope=0.05, xscale=1.6, yscale=1.4, timing/font=\sffamily\small, timing/rowdist=3.3ex]
  clk   & 6{1L 1H} 3S 6{1L 1H} \\
  cnt   & 2D{0} 2D{1} 2D{2} 2D{\ldots} 2D{H$-$2} 2D{H$-$1} 3S 2D{0} 2D{1} 2D{2} 2D{\ldots} 2D{H$-$2} 2D{H$-$1} \\
  state & 12L 3S 12H \\
\extracode
  \node[lbl, text=FpgaOrange] at (13.5,0.7) {$\approx$0.5\,s};
\end{tikztimingtable}
\caption{Timing diagram of the blinker (H = HALF)}
\end{figure}

\begin{figure}[H]
\centering
\begin{tikzpicture}
\begin{axis}[
  width=0.85\textwidth, height=4.5cm, xlabel={Time, s}, ylabel={led},
  ymin=-0.2, ymax=1.3, xmin=0, xmax=3, ytick={0,1}, yticklabels={on (0), off (1)},
  label style={font=\sffamily\small}, tick label style={font=\sffamily\small},
  grid=major, grid style={FpgaGray!20}, axis lines*=left,
]
  \addplot[very thick, FpgaBlue, const plot] coordinates {(0,0) (0.5,1) (1,0) (1.5,1) (2,0) (2.5,1) (3,1)};
\end{axis}
\end{tikzpicture}
\caption{The signal on the LED pin: period 1\,s, frequency 1\,Hz}
\end{figure}

\begin{taskbox}
\begin{enumerate}
  \item Make the LED blink 5 times per second.
  \item Make \code{led[0]} blink at 1\,Hz, \code{led[1]} at 2\,Hz,
        \code{led[2]} at 4\,Hz\ldots\ Hint: use different bits of a single
        counter. Which bit of the 25-bit counter changes at a frequency of about 1\,Hz?
\end{enumerate}
\end{taskbox}

\begin{diybox}[The ``armed'' blinker]
\textbf{Task.} The LED of the armed mode must blink once per second while the panel is armed. Choose the divider and check the period.

\textbf{In Logra.} The ``From an example'' card---``27 MHz blinker'', or your own \code{blink.v}. Code tab $\to$ ``Synthesize''; in Waveforms reduce the divider constant (the \code{CLK\_HZ} parameter), otherwise you would have to simulate millions of clock cycles.

\textbf{Expected result.} The LED toggles every $13\,500\,000$ clock cycles (0.5\,s), that is, it blinks with a period of 1\,s and a frequency of 1\,Hz; in the model with \code{CLK\_HZ = 20} it toggles every 10 clock cycles.

\textbf{How to check.} Waveforms: set the number of cycles to 25 and watch \code{cnt} and \code{led}; \code{cnt} reaches \code{HALF - 1} and resets. On the Chip tab, after ``Place and build'', make sure that $f_{max}$ is much higher than 27\,MHz.

\stars{1} repeat: a period of 1\,s from the diagram. \stars{2} change: 5 blinks per second; then find which bit of the counter toggles at a frequency of about 1\,Hz. \stars{3} project: the ``Security Panel'' stage---a mode indicator; acceptance criterion: the LED blinks with a period of 1\,s only when \code{armed = 1}, otherwise it is off.
\end{diybox}

% =============================================================================
\section{Project 3. A Running Light}

\textbf{Goal:} to study the shift register and the enable signal (\code{tick}).

Six flip-flops are connected in a ring: on each \code{tick} pulse the value of
each flip-flop moves to its neighbor, and the value of the last one goes to the first.
If the ring initially holds a single one, it will ``run'' around the circle.

\begin{figure}[H]
\centering
\begin{tikzpicture}[ffn/.style={draw=FpgaGreen, fill=green!8, thick, minimum width=13mm, minimum height=11mm, font=\sffamily\small, align=center}]
  \foreach \i in {0,...,5} {
    \node[ffn] (r\i) at (\i*2.2,0) {D\enskip Q\\[-1mm]\footnotesize ring[\i]};
    \draw[wire] (r\i.north) -- ++(0,0.6) node[above, font=\sffamily\footnotesize] {led[\i]};
  }
  \foreach \i [evaluate=\i as \j using int(\i+1)] in {0,...,4}
    \draw[arr] (r\i.east) -- (r\j.west);
  \draw[arr] (r5.east) -- ++(0.4,0) |- ($(r0.south)+(0,-0.7)$) -| ($(r0.west)+(-0.4,0)$) -- (r0.west);
  \draw[wire, FpgaBlue] ($(r0.south)+(0,-1.3)$) node[left, font=\sffamily\footnotesize] {clk, tick} -- ($(r5.south)+(0,-1.3)$);
  \foreach \i in {0,...,5} \draw[wire, FpgaBlue] (r\i.south) -- ++(0,-1.3);
\end{tikzpicture}
\caption{A ring shift register of six D flip-flops}
\end{figure}

\vfile{\codepath/03_running_lights/running_lights.v}{code/03\_running\_lights/running\_lights.v}

The pin file is the same as for the blinker: the clock and the LEDs.

\vfile[c]{\codepath/03_running_lights/running_lights.cst}{code/03\_running\_lights/running\_lights.cst}

\begin{figure}[H]
\centering
\begin{tikztimingtable}[timing/slope=0.1, xscale=1.35, yscale=1.35, timing/font=\sffamily\small, timing/rowdist=3.1ex]
  tick    & 2L 1H 3L 1H 3L 1H 3L 1H 3L 1H 3L 1H 3L 1H 1L \\
  ring[0] & 3H 20L 4H 1L \\
  ring[1] & 3L 4H 21L \\
  ring[2] & 7L 4H 17L \\
  ring[3] & 11L 4H 13L \\
  ring[4] & 15L 4H 9L \\
  ring[5] & 19L 4H 5L \\
\end{tikztimingtable}
\caption{The one ``moves'' to the neighboring bit after each \code{tick} pulse
(the time scale is schematic)}
\end{figure}

\begin{infobox}[Why tick and not a separate slow clock?]
You could feed the flip-flops a homemade slow clock taken straight from
the divider. But this is not done in FPGAs: such a signal has no access to the
global clock network, and the circuit may work unreliably. The correct approach is that
\textbf{all flip-flops are clocked by the single \code{clk}}, and rare events
are expressed by an \term{enable signal} one clock cycle long.
\end{infobox}

\begin{taskbox}
Make a light that runs back and forth (``Knight Rider''). Hint:
you need an additional direction register \code{dir}.
\end{taskbox}

\begin{diybox}[A running light in both directions]
\textbf{Task.} The light must run to the last LED and back, and in the alarm state it must also run faster.

\textbf{In Logra.} The ``From an example'' card---``Running light''; Code tab. Add a register \code{dir} that changes its value when the one reaches an edge.

\textbf{Expected result.} The one goes $0\to5$, then $5\to0$, and so on in a loop; at each step exactly one bit equals 1.

\textbf{How to check.} Waveforms: watch \code{ring} and \code{tick}, as in the figure for this project; the one should ``bounce'' off the edges. RTL shows six flip-flops in the ring.

\stars{1} repeat: the example unchanged, shifting on the \code{tick} signal. \stars{2} change: direction and speed controlled by buttons (Task~6 of Chapter~\ref{ch:tasks}). \stars{3} project: the ``Security Panel'' stage---an alarm indicator; acceptance criterion: when \code{alarm = 1} the light runs twice as fast, and outside the alarm the LEDs are off.
\end{diybox}

% =============================================================================
\section{Project 4. The Button: Synchronization, Bounce, Edge}

\textbf{Goal:} to learn to work with buttons correctly. The task is simple:
count presses of S1 and show the number in binary on the LEDs. But
there are three obstacles in the way.

\subsection{Obstacle 1: Contact Bounce}

When pressed, the mechanical contacts of a button ``jump'' for several milliseconds, and
instead of one $0\to1$ transition you get a whole series. A circuit running
at 27\,MHz will count every ``jump'' as a separate press.

\begin{figure}[H]
\centering
\begin{tikzpicture}
\begin{axis}[
  width=0.9\textwidth, height=4.6cm, xlabel={Time, ms}, ylabel={Level},
  xmin=0, xmax=30, ymin=-0.2, ymax=1.5, ytick={0,1},
  label style={font=\sffamily\small}, tick label style={font=\sffamily\small},
  legend style={font=\sffamily\small, at={(0.5,1.08)}, anchor=south, legend columns=2, draw=none},
  axis lines*=left,
]
  \addplot[thick, FpgaRed, const plot] coordinates {
    (0,0) (5,1) (5.4,0) (5.9,1) (6.2,0) (6.9,1) (7.1,0) (7.8,1) (8.4,0) (8.6,1)
    (22,1) (22.3,0) (22.7,1) (23.2,0) (23.5,1) (24.0,0) (30,0)};
  \addlegendentry{button (raw signal)}
  \addplot[very thick, FpgaGreen, const plot] coordinates {(0,-0.05) (18.6,-0.05) (18.6,1.1) (30,1.1)};
  \addlegendentry{after debouncing}
  \draw[FpgaOrange, <->, thick] (axis cs:8.6,1.25) -- (axis cs:18.6,1.25) node[midway, above, font=\sffamily\small] {10 ms stable};
\end{axis}
\end{tikzpicture}
\caption{Contact bounce and its suppression: a new value is accepted only
if it does not change for 10 ms}
\end{figure}

\subsection{Obstacle 2: Metastability}

A button is pressed at a random moment that has nothing to do with the \code{clk} clock.
If the input of a flip-flop changes exactly at the moment of the edge, the flip-flop may
``hang'' for a while between 0 and 1 (the \term{metastable state}). The protection
is a \term{synchronizer} made of two flip-flops in series: even
if the first one ``hesitates'', by the second edge it will almost certainly have settled.

\begin{figure}[H]
\centering
\begin{tikzpicture}[ffn/.style={draw=FpgaGreen, fill=green!8, thick, minimum width=13mm, minimum height=11mm, font=\sffamily\small}]
  \node[font=\sffamily\small] (in) at (0,0) {btn};
  \node[ffn] (f1) at (2.2,0) {D\enskip Q};
  \node[ffn] (f2) at (4.6,0) {D\enskip Q};
  \node[block] (db) at (7.6,0) {debounce\\\footnotesize 10 ms};
  \node[block] (ed) at (10.8,0) {edge\\detector};
  \draw[arr] (in) -- (f1); \draw[arr] (f1) -- node[above, lbl] {sync0} (f2);
  \draw[arr] (f2) -- node[above, lbl] {sync1} (db); \draw[arr] (db) -- node[above, lbl] {s1} (ed);
  \draw[arr] (ed) -- ++(1.6,0) node[right, font=\sffamily\small] {press};
  \draw[wire, FpgaBlue] (1.2,-1.2) node[left, font=\sffamily\footnotesize] {clk} -- (10.8,-1.2);
  \foreach \n in {f1,f2,db,ed} \draw[wire, FpgaBlue] (\n.south) -- (\n.south |- 0,-1.2);
\end{tikzpicture}
\caption{The button signal processing chain}
\end{figure}

\subsection{Obstacle 3: A Press Lasts Millions of Clock Cycles}

If you write \code{if (s1) count <= count + 1;}, the counter will add
one \emph{every clock cycle} while the button is held---millions of times.
You need a \term{rising-edge detector}: compare the current value with the
value delayed by one clock cycle. The pulse \code{press = s1 \& \textasciitilde{}s1\_prev} lasts
exactly one clock cycle.

\begin{figure}[H]
\centering
\begin{tikztimingtable}[timing/slope=0.1, xscale=2.4, yscale=1.4, timing/font=\sffamily\small, timing/rowdist=3.2ex]
  clk      & 8{1L 1H} \\
  s1       & 3L 10H 3L \\
  s1\_prev & 5L 10H 1L \\
  press    & 3L 2H 11L \\
  count    & 5D{5} 11D{6} \\
\end{tikztimingtable}
\caption{The edge detector: \code{press} equals 1 for only one clock cycle after the press}
\end{figure}

\vfile{\codepath/04_button_edge/button_counter.v}{code/04\_button\_edge/button\_counter.v}

The pin file: the clock, both buttons and the LEDs. The \code{PULL\_MODE=DOWN} pull-down
does not save you from bounce---it only sets the level of a released button.
Bounce is handled by the debounce circuit in the code.

\vfile[c]{\codepath/04_button_edge/button_counter.cst}{code/04\_button\_edge/button\_counter.cst}

\begin{taskbox}
\begin{enumerate}
  \item Make S2 subtract one from the counter instead of resetting it. For this,
        S2 also needs debouncing and an edge detector. It is convenient to move them
        into a separate \code{debounce} module and create two instances.
  \item What happens if you remove the debouncing? Try it and count how many
        extra ``presses'' accumulate.
\end{enumerate}
\end{taskbox}

\begin{diybox}[Debouncing and counting presses]
\textbf{Task.} The ``arm the system'' button must produce exactly one transition per press, no matter how the contacts bounce.

\textbf{In Logra.} The ``From an example'' card: ``Button debounce'' and ``One press $	o$ toggle''. Code tab, then Waveforms.

\textbf{Expected result.} In ``Button debounce'' the \code{synced} signal repeats the noise at the input, while \code{stable} changes only after several consistent samples. In ``One press $	o$ toggle'' a long press gives one \code{press} pulse and one toggle.

\textbf{How to check.} Waveforms: apply the series 1-0-1-0-1 to the button and then a steady one; compare \code{synced} and \code{stable}. Remove the debouncing from \code{button\_counter.v} and count the ``extra'' presses.

\stars{1} repeat: run both examples and find the clock cycle where \code{press = 1}. \stars{2} find the bug: remove the two-flip-flop synchronizer and explain why this is dangerous on a board (metastability, Chapter~\ref{ch:clocks}). \stars{3} project: the ``Security Panel'' stage---the arming button; acceptance criterion: on the diagram, with bounce at the input, the panel state changes exactly once.
\end{diybox}

% =============================================================================
\section{Project 5. A Traffic Light: A Finite-State Machine}

\textbf{Goal:} to move onto the board the finite-state machine that you studied in
Chapter~\ref{ch:sync} and to see it work on the LEDs. The principle is the same: a state
register, transition logic, output logic. The only new things here are the timer and the
connection to real pins.

A machine is in one of a finite number of \emph{states} and, under
certain conditions, \emph{moves} to others. A traffic light is a classic
example:

\begin{figure}[H]
\centering
\begin{tikzpicture}[
  st/.style={circle, draw=FpgaDark, very thick, minimum size=22mm, font=\sffamily\small, align=center},
  tr/.style={-{Stealth[length=3mm]}, very thick, FpgaDark},
]
  \node[st, fill=red!20]    (r)  at (0,2.6)  {RED\\\footnotesize 4 s};
  \node[st, fill=orange!25] (ry) at (5.2,2.6) {RED\_YL\\\footnotesize 1 s};
  \node[st, fill=green!20]  (g)  at (5.2,-1.6) {GREEN\\\footnotesize 4 s};
  \node[st, fill=yellow!30] (y)  at (0,-1.6)  {YELLOW\\\footnotesize 1 s};
  \draw[tr] (r)  -- node[above, lbl] {timeout} (ry);
  \draw[tr] (ry) -- node[left, lbl] {timeout} (g);
  \draw[tr] (g)  -- node[below, lbl] {timeout} (y);
  \draw[tr] (y)  -- node[left, lbl] {timeout} (r);
  \draw[tr, FpgaRed, dashed] (-2.6,4.2) node[left, lbl, text=FpgaRed] {btn[0] (reset)} -- (r);
  % Lamps
  \foreach \s/\dx/\dy/\cr/\cy/\cg in {r/-1.7/0/red/gray!30/gray!30, ry/1.7/0.5/red/yellow/gray!30, g/1.7/0.5/gray!30/gray!30/green!70!black, y/-1.7/0/gray!30/yellow/gray!30} {
    \begin{scope}[shift={($(\s)+(\dx,\dy)$)}]
    \fill[ChipBlack, rounded corners=1pt] (-0.22,-1.0) rectangle (0.22,0.35);
    \fill[\cr] (0,0.15) circle (0.13); \fill[\cy] (0,-0.3) circle (0.13); \fill[\cg] (0,-0.75) circle (0.13);
    \end{scope}
  }
\end{tikzpicture}
\caption{The state transition graph of the ``Traffic light'' machine}
\end{figure}

In Verilog a machine is described in three parts:
\begin{enumerate}
  \item the \textbf{state register} (\code{state})---flip-flops that store
        the number of the current state;
  \item the \textbf{transition logic}---determines the next state;
  \item the \textbf{output logic}---what to light in each state.
\end{enumerate}

\begin{figure}[H]
\centering
\begin{tikzpicture}[node distance=12mm]
  \node[block, minimum width=30mm, minimum height=14mm] (next) {Transition\\logic};
  \node[gblock, right=of next, minimum height=14mm] (st) {Register\\state};
  \node[block, right=of st, minimum width=30mm, minimum height=14mm] (out) {Output\\logic};
  \draw[arr] (next) -- (st); \draw[arr] (st) -- (out);
  \draw[arr] (out.east) -- ++(1,0) node[right, font=\sffamily\small] {led};
  \draw[arr] ($(st.east)+(0.5,0)$) -- ++(0,-1.2) -| (next.south);
  \draw[arr] ($(next.west)+(-1.2,0)$) node[left, font=\sffamily\small] {timer, btn} -- (next.west);
  \draw[wire, FpgaBlue] (st.north) -- ++(0,0.5) node[above, font=\sffamily\footnotesize] {clk};
\end{tikzpicture}
\caption{Structure of a Moore machine: the output depends only on the state}
\end{figure}

\vfile{\codepath/05_traffic_light/traffic_light.v}{code/05\_traffic\_light/traffic\_light.v}

The pin file: the clock, the buttons (S1 is used as reset) and the LEDs. The
\code{btn} port in the module is two bits wide, so both bits are described even though \code{btn[1]}
is not used: otherwise the router will report that the port is not described.

\vfile[c]{\codepath/05_traffic_light/traffic_light.cst}{code/05\_traffic\_light/traffic\_light.cst}

\begin{figure}[H]
\centering
\begin{tikztimingtable}[timing/slope=0.05, xscale=1.2, yscale=1.4, timing/font=\sffamily\small, timing/rowdist=3.3ex]
  state   & 8D{RED} 2D{R+Y} 8D{GREEN} 2D{YEL} 8D{RED} 2D{R+Y} \\
  red     & 10H 10L 10H \\
  yellow  & 8L 2H 8L 2H 8L 2H \\
  green   & 10L 8H 12L \\
\extracode
  \draw[FpgaGray] (0,-6.2) -- (30,-6.2);
  \foreach \t/\l in {0/0, 8/4, 10/5, 18/9, 20/10, 28/14, 30/15} {
    \draw[FpgaGray] (\t,-6.1) -- (\t,-6.3) node[below, font=\sffamily\scriptsize] {\l\,s};
  }
\end{tikztimingtable}
\caption{The traffic light over time (1 = lamp on)}
\end{figure}

\begin{taskbox}
Add a ``flashing green'' state (3 times for 0.5\,s each) between GREEN and
YELLOW. Draw the new transition graph before you write the code!
\end{taskbox}

\begin{diybox}[Exit delay in the security system]
\textbf{Task.} After the system is armed, the owner has 3\,s to leave: during this time the sensors are ignored. Add an ``exit'' state to the panel's machine.

\textbf{In Logra.} The ``From an example'' card---``Traffic light state machine''---as a model (\code{traffic\_light.v}); Code tab. Draw the graph OFF $\to$ EXIT $\to$ ARMED $\to$ ALARM, then write the \code{case}.

\textbf{Expected result.} After \code{arm} the machine stays in EXIT for exactly 3\,s (30 clock cycles in the reduced model), then moves to ARMED; a sensor trigger in EXIT does not affect the state.

\textbf{How to check.} Waveforms: the \code{state} signal goes through $0\to 1\to 2\to 3$; on the diagram, a sensor raised in EXIT does not cause a transition to ALARM. RTL shows the state register and the timer counter.

\stars{1} repeat: go through red, red-yellow, green, yellow on the diagram. \stars{2} change: add ``flashing green'' (the task at the end of the project). \stars{3} project: the ``Security Panel'' stage---exit delay; acceptance criterion: a graph of four states is drawn, and the diagram has no transitions that are not on the graph.
\end{diybox}

% =============================================================================
\section{Project 6. A ``Breathing'' LED: PWM}

\textbf{Goal:} to learn to control brightness. A digital pin has only two
states: 0 and 1. But if you switch an LED very fast, the eye sees the
\emph{average} brightness. This is \term{pulse-width modulation} (PWM):
the frequency is constant, while the fraction of the time when the signal equals 1 changes---the
\term{duty cycle} $D$:
\[
  D = \frac{t_\text{on}}{T}\cdot 100\,\%, \qquad
  U_\text{avg} = D \cdot U_\text{supply}.
\]

\begin{figure}[H]
\centering
\begin{tikzpicture}
\begin{groupplot}[
  group style={group size=1 by 3, vertical sep=13mm, xticklabels at=edge bottom},
  width=0.85\textwidth, height=3cm, xmin=0, xmax=4, ymin=-0.15, ymax=1.3,
  ytick={0,1}, axis lines*=left, tick label style={font=\sffamily\small},
  label style={font=\sffamily\small},
]
\nextgroupplot[title={$D=25\,\%$ --- dim}, title style={font=\sffamily\small}]
  \addplot[very thick, FpgaBlue, const plot] coordinates {(0,1) (0.25,0) (1,1) (1.25,0) (2,1) (2.25,0) (3,1) (3.25,0) (4,0)};
  \addplot[thick, dashed, FpgaOrange, domain=0:4] {0.25};
\nextgroupplot[title={$D=50\,\%$}, title style={font=\sffamily\small}]
  \addplot[very thick, FpgaBlue, const plot] coordinates {(0,1) (0.5,0) (1,1) (1.5,0) (2,1) (2.5,0) (3,1) (3.5,0) (4,0)};
  \addplot[thick, dashed, FpgaOrange, domain=0:4] {0.5};
\nextgroupplot[title={$D=90\,\%$ --- bright}, title style={font=\sffamily\small}, xlabel={Time, PWM periods}]
  \addplot[very thick, FpgaBlue, const plot] coordinates {(0,1) (0.9,0) (1,1) (1.9,0) (2,1) (2.9,0) (3,1) (3.9,0) (4,0)};
  \addplot[thick, dashed, FpgaOrange, domain=0:4] {0.9};
\end{groupplot}
\end{tikzpicture}
\caption{PWM with different duty cycles. The dashed line is the average value
that the eye ``sees''}
\end{figure}

How do you get PWM on an FPGA? You need a free-running counter and a comparator:
while the counter is less than a given number \code{duty}, the output equals 1.

\begin{figure}[H]
\centering
\begin{tikzpicture}
\begin{axis}[
  width=0.9\textwidth, height=5.2cm, xmin=0, xmax=768, ymin=0, ymax=300,
  xlabel={Clock cycles}, ylabel={Value},
  label style={font=\sffamily\small}, tick label style={font=\sffamily\small},
  legend style={font=\sffamily\small, at={(0.5,1.02)}, anchor=south, legend columns=3, draw=none},
  axis lines*=left,
]
  \addplot[thick, FpgaBlue] coordinates {(0,0) (255,255) (256,0) (511,255) (512,0) (767,255)};
  \addlegendentry{pwm\_cnt}
  \addplot[thick, dashed, FpgaRed] coordinates {(0,96) (768,96)};
  \addlegendentry{duty = 96}
  \addplot[very thick, FpgaGreen, const plot] coordinates {(0,285) (96,265) (256,285) (352,265) (512,285) (608,265) (768,265)};
  \addlegendentry{pwm (output)}
\end{axis}
\end{tikzpicture}
\caption{A sawtooth counter and a comparator: the \code{pwm} output equals 1 while
\code{pwm\_cnt < duty}}
\end{figure}

To make the LED ``breathe'', we will slowly increase \code{duty} from 0 to 255
and then decrease it back.

\begin{figure}[H]
\centering
\begin{tikzpicture}
\begin{axis}[
  width=0.85\textwidth, height=4.6cm, xmin=0, xmax=6, ymin=0, ymax=270,
  xlabel={Time, s}, ylabel={duty},
  label style={font=\sffamily\small}, tick label style={font=\sffamily\small},
  axis lines*=left, grid=major, grid style={FpgaGray!20},
]
  \addplot[very thick, FpgaPurple] coordinates {(0,0) (1,255) (2,0) (3,255) (4,0) (5,255) (6,0)};
\end{axis}
\end{tikzpicture}
\caption{Change of the duty cycle: an ``inhale'' in 1\,s and an ``exhale'' in 1\,s}
\end{figure}

\vfile{\codepath/06_pwm_breath/pwm_breath.v}{code/06\_pwm\_breath/pwm\_breath.v}

The pin file: the clock and the LEDs. \code{DRIVE=8} limits the output current;
PWM brightness does not depend on it---it is set by the fraction of time the
LED is on.

\vfile[c]{\codepath/06_pwm_breath/pwm_breath.cst}{code/06\_pwm\_breath/pwm\_breath.cst}

\begin{infobox}[Why the LED seems ``uneven'']
The eye perceives brightness nonlinearly: the difference between $D=5\,\%$ and $10\,\%$ is noticeable
much more than between $90\,\%$ and $95\,\%$. To make the ``breathing'' look
smooth, real devices use \term{gamma correction}: a table
that converts linear brightness into the required \code{duty} value. Such a table
is conveniently stored in BSRAM block memory.
\end{infobox}

\begin{taskbox}
\begin{enumerate}
  \item Make the six LEDs ``breathe'' with a phase shift
        (a ``wave'').
  \item Add control: button S1 increases the brightness by 1/8, S2
        decreases it (use the debounce module from project 4).
\end{enumerate}
\end{taskbox}

\begin{diybox}[Breathing of the armed indicator]
\textbf{Task.} The ``armed'' indicator must smoothly ``breathe'', and the siren volume must be adjustable with the number \code{duty}.

\textbf{In Logra.} The ``From an example'' card---``Brightness with PWM'' and ``Smooth LED breathing''. Code tab, then Waveforms with a large number of cycles.

\textbf{Expected result.} With \code{duty = 96} out of 256 clock cycles the output equals 1 for the first 96 cycles of the period (37.5\,\%); as \code{duty} grows from 0 to 255 the brightness rises, then falls.

\textbf{How to check.} Waveforms: signals \code{pwm\_cnt}, \code{duty}, \code{pwm}; compare the fraction of ones with the formula $D=t_\text{on}/T$. Brightness on the board is judged by eye: the larger \code{duty}, the brighter.

\stars{1} repeat: measure $D$ for three values of \code{duty}. \stars{2} change: a ``wave'' on six LEDs with a phase shift. \stars{3} project: the ``Security Panel'' stage---indication; acceptance criterion: in the ARMED mode the indicator breathes with a period of about 2\,s, and in OFF it is off.
\end{diybox}

% =============================================================================
\section{Project 7. Decoder, Multiplexer and Demultiplexer}

\textbf{Goal:} to put together the blocks from Chapters~\ref{ch:coders} and~\ref{ch:mux} in an FPGA and
to see them work on the LEDs.

Buttons S1 and S2 form a two-bit address $s$ (from 0 to 3). A divider counter
produces four signals of different frequencies---these are its most significant bits
(Fig.~\ref{fig:count_t}).
\begin{itemize}
  \item The \textbf{4:1 multiplexer} uses the address to select one of the four signals
        and sends it to \code{led[4]}: the higher the number, the faster the blinking.
  \item The \textbf{2$\to$4 decoder} turns the address into a one-hot code.
  \item The \textbf{1:4 demultiplexer} is the same decoder whose
        outputs are ``enabled'' by a data signal. The blinking reaches only
        the LED \code{led[s]}.
\end{itemize}

\begin{figure}[H]
\centering
\begin{tikzpicture}
  \node[hblock, minimum height=24mm, minimum width=18mm] (cnt) at (0,0) {Counter\\\footnotesize 25 bits};
  \draw[thick, fill=orange!15] (2.6,1.4) -- (3.6,0.9) -- (3.6,-0.9) -- (2.6,-1.4) -- cycle;
  \node[font=\sffamily\footnotesize] at (3.1,0) {MUX};
  \foreach \i in {0,...,3} \draw[wire] (cnt.east |- 0,0.9-\i*0.6) -- node[above, font=\footnotesize, pos=0.4] {cnt[\the\numexpr21+\i\relax]} (2.6,0.9-\i*0.6);
  \draw[wire] (3.6,0) -- (4.4,0) node[right, font=\sffamily\small] {led[4]};
  \draw[thick, fill=FpgaLight] (6.4,-5.4) rectangle (7.4,-0.6);
  \node[font=\sffamily\small\bfseries] at (6.9,-3.0) {DC};
  \foreach \i in {0,...,3} {
    \node[gAND, anchor=input 1] (g\i) at (9.2,-1.0-\i*1.2) {};
    \draw[wire] (7.4,-1.0-\i*1.2) -- (g\i.input 1);
    \draw[wire] (g\i.output) -- ++(0.4,0) node[right, font=\sffamily\small] {led[\i]};
  }
  \draw[wire, FpgaOrange] (cnt.south) -- ++(0,-5.0) node[below, font=\sffamily\footnotesize, text=FpgaDark] {cnt[23]} -| (8.1,-1.45);
  \foreach \i in {0,...,3} { \draw[wire, FpgaOrange] (8.1,0 |- g\i.input 2) -- (g\i.input 2); \fill[FpgaOrange] (8.1,0 |- g\i.input 2) circle (1.2pt); }
  \draw[wire, FpgaOrange] (8.1,-1.45) -- (8.1,0 |- g3.input 2);
  \draw[wire, FpgaBlue] (4.9,-3.0) node[left, text=FpgaDark] {s = btn} -- (6.4,-3.0);
  \draw[wire, FpgaBlue] (5.4,-3.0) |- (3.1,-2.0) -- (3.1,-1.15); \fill[FpgaBlue] (5.4,-3.0) circle (1.3pt);
  \begin{scope}[on background layer]
    \node[draw=FpgaGray, dashed, rounded corners, fit=(g0)(g3), inner sep=6mm, label={[lbl]above:demultiplexer}] {};
  \end{scope}
\end{tikzpicture}
\caption{The circuit of project 7. Demultiplexer = decoder + AND gates}
\end{figure}

\vfile{\codepath/07_mux_demux/mux_demux.v}{code/07\_mux\_demux/mux\_demux.v}

The pin file: the clock, both buttons (they form the address $s$) and the LEDs.

\vfile[c]{\codepath/07_mux_demux/mux_demux.cst}{code/07\_mux\_demux/mux\_demux.cst}

Note that the same block can be described in different ways:
the multiplexer here is written as a tree of \code{?:} operators (as in
Fig.~\ref{fig:muxtree}), and the decoder as a \code{case}. The shortest way to write a
multiplexer is \code{assign y = x[s];}---the synthesizer will produce
the same circuit.

\begin{infobox}[What we got after synthesis]
Yosys placed the multiplexer not in ordinary LUTs but in the special
multiplexers \code{MUX2\_LUT5} and \code{MUX2\_LUT6}, which sit in every
logic cell next to the LUT. The whole project took 8 LUTs and 25 flip-flops
(the flip-flops are the counter).
\end{infobox}

\begin{taskbox}
\begin{enumerate}
  \item Replace the multiplexer with the line \code{assign y = x[s];} and compare
        the synthesis reports.
  \item Add a priority encoder (Chapter~\ref{ch:coders}): output on
        \code{led[5]} the flag $V$ ``at least one button is pressed''.
  \item Show the number $s$ on a seven-segment display connected to the
        GPIO pins, using the decoder from Table~\ref{tab:7seg}.
\end{enumerate}
\end{taskbox}

\begin{diybox}[Security zones: multiplexer and demultiplexer]
\textbf{Task.} The security panel monitors four zones. Buttons S1, S2 select the zone number $s$; the multiplexer shows the state of the selected zone on one LED, and the demultiplexer lights the lamp of exactly that zone.

\textbf{In Logra.} The ``From an example'' card: ``Choosing one of four words'', ``Demultiplexer 1 $	o$ 4'', ``2 $	o$ 4 decoder''; Code tab---\code{mux\_demux.v} from the book. The Pins tab holds the assignment of buttons and LEDs.

\textbf{Expected result.} \code{led[4]} blinks at a frequency that depends on $s$ (the larger $s$, the faster). The blinking reaches only \code{led[s]}; the other three are off at that moment.

\textbf{How to check.} Waveforms: sweep $s=0\ldots3$ while the counter is running. RTL: the multiplexer, the decoder and the demultiplexer's AND gates. On the Gowin tab find the \code{MUX2\_LUT5} and \code{MUX2\_LUT6} cells.

\stars{1} repeat: check all four addresses. \stars{2} change: replace the multiplexer with \code{assign y = x[s];} and compare the synthesis reports. \stars{3} project: the ``Security Panel'' stage---zone selection; acceptance criterion: when the sensor of zone $z$ triggers, the lamp of zone $z$ is on and the others are off, and the multiplexer outputs that zone's sensor to the indicator by zone number.
\end{diybox}
